\chapter{Decomposing the Spread}\label{ch:mx:decomposing-the-spread}

A buy of 1,000 shares pays two cents over the mid. A minute later the mid has moved one and a half cents the way the trade went: the liquidity
provider kept half a cent, and the rest was what the trade told the market. Chapter~4 explained why a spread exists; this chapter measures what each
part of it pays for. The measurement is harder than it looks: the move after a trade mixes what the trade revealed, the trades that followed it and
the market's own noise, and each method untangles them under its own assumptions. We test every method on \texttt{firm.tape}, the simulated market of
One Quant Book 7, where the efficient price is known and the true adverse selection of every trade can be computed.

\section{Quoted, effective and realised spreads}

\begin{definition}[Quoted spread]\label{def:mx:decomposing-the-spread:quoted}
The \emph{quoted spread}\index{quoted spread} at time $t$ is $a_t-b_t$, the best ask minus the best bid; the relative quoted spread divides it by the
mid $m_t=\tfrac12(a_t+b_t)$. Averaged over a period, it is weighted by the time each quote was in force.
\end{definition}

The \omterm{def:mx:decomposing-the-spread:quoted}{quoted spread} is what a small order would pay; the effective spread (One Quant Book 1, chapter~10) is what trades did pay. With $\varepsilon_t=+1$
for a buyer-initiated trade and $-1$ for a seller-initiated one (the trade sign, One Quant Book 7, chapter~9), a trade at price $p_t$ has effective
half-spread $\varepsilon_t(p_t-m_t)$, measured against the mid just before it. It is below the quoted half-spread when the trade received price
improvement, above it when the order walked the book. The realised spread asks what the liquidity provider kept once the price had moved.

\begin{definition}[Trade price impact]\label{def:mx:decomposing-the-spread:impact}
The \emph{trade price impact}\index{trade price impact} of a trade at horizon $h$ is $\varepsilon_t(m_{t+h}-m_t)$: the move of the mid in the trade's
direction over the next $h$. The realised half-spread is $\varepsilon_t(p_t-m_{t+h})$, what the provider would earn by closing her position at the later
mid, and trade by trade
\[
\underbrace{\varepsilon_t(p_t-m_t)}_{\text{effective}}=\underbrace{\varepsilon_t(p_t-m_{t+h})}_{\text{realised}}+\underbrace{\varepsilon_t(m_{t+h}-m_t)}_{\text{impact}}.
\]
\end{definition}

\begin{omfigure}
\begin{tikzpicture}[font=\footnotesize,x=1cm,y=1.1cm]
\draw[->,gray] (0,0) -- (9.5,0) node[right] {time};
\draw[omDef,very thick] (0.3,1) -- (3,1) node[pos=0.1,above] {mid $m_t$};
\draw[omDef,very thick] (3,1) -- (4.2,1.5) -- (5.4,1.9) -- (7,2.5) -- (9.2,2.5);
\node[omDef,right] at (9.2,2.5) {mid $m_{t+h}$};
\fill[omThm] (3,3) circle (2.5pt) node[above] {buy at $p_t$};
\draw[dashed,gray] (3,0) -- (3,3); \node[below] at (3,0) {$t$};
\draw[dashed,gray] (7.8,0) -- (7.8,2.5); \node[below] at (7.8,0) {$t+h$};
\draw[dotted] (3,3) -- (9.2,3);
\draw[decorate,decoration={brace,amplitude=4pt,mirror}] (2.8,1) -- (2.8,3) node[midway,left=5pt,align=right] {effective\\ half-spread};
\draw[decorate,decoration={brace,amplitude=4pt}] (8.1,3) -- (8.1,2.5) node[midway,right=5pt] {realised};
\draw[decorate,decoration={brace,amplitude=4pt}] (8.1,2.5) -- (8.1,1) node[midway,right=5pt] {impact};
\draw[dotted] (3,1) -- (8.1,1);
\end{tikzpicture}

\omcaption{One buy, split in two. The effective half-spread is what the buyer paid over the mid; after $h$, the part the mid has moved in the trade's
direction is the impact, the part left to the liquidity provider is the realised spread.}\label{fig:mx:decomposing-the-spread:split}
\end{omfigure}

\omcode{code/firm/spreaddecomp/firm_spreaddecomp.py}{41}{48}{Share-weighted effective, realised and impact half-spreads at one horizon.}\label{lst:mx:decomposing-the-spread:decompose}

The simulated session has 6.5 hours, 21\,993 trades of 118 shares on average, and a time-weighted \omterm{def:mx:decomposing-the-spread:quoted}{quoted spread} of 1.04 ticks. The trades' effective
half-spread, weighted by shares as regulatory reports weight it, is 0.535 ticks.

\begin{dated}{2026-09}{Realised spreads in Rule 605 reports}\label{dat:mx:decomposing-the-spread:rule605}
The SEC's amendments to Rule 605 of Regulation NMS (Release 34-99679, adopted March 6, 2024, published April 15, 2024) require execution-quality
reports to give realised spreads at five horizons, 50 milliseconds, 1 second, 15 seconds, 1 minute and 5 minutes, instead of the single 5-minute
horizon of the original rule, and to measure times in milliseconds or finer. The compliance date, first December 14, 2025, was extended on
September 30, 2025 to August 1, 2026.
\end{dated}

\section{The price impact of a trade}

\begin{definition}[Adverse-selection component]\label{def:mx:decomposing-the-spread:adverse}
The \emph{adverse-selection component}\index{adverse-selection component} of the effective spread is the part that compensates liquidity providers for
trading with counterparties who know more: for a trade at $t$, the expected value of $\varepsilon_t(P^\ast_t-m_t)$, where $P^\ast_t$ is the efficient price,
the expectation of the asset's value given all information, including what the aggressor knows and the quotes do not yet show.
\end{definition}

\begin{definition}[Spread decomposition]\label{def:mx:decomposing-the-spread:decomposition}
A \emph{spread decomposition}\index{spread decomposition} splits the effective spread into an \omterm{def:mx:decomposing-the-spread:adverse}{adverse-selection component} and the rest, the
order-processing and \omterm{def:mx:why-there-is-a-spread:inventory}{inventory-holding costs} of chapter~4, which the provider keeps once the price has settled. Methods differ in what they observe
(trades only, trades and quotes, daily bars) and in the model that lets them tell the parts apart.
\end{definition}

In a real market $P^\ast$ is not observed, and the realised-spread method stands in for it with a later mid. It works when the mid catches up with the
efficient price and nothing else moves it in the trade's direction.

\begin{proposition}[What the impact converges to]\label{prop:mx:decomposing-the-spread:converge}
Let $P^\ast$ be a martingale, let $\varepsilon_t$ and $m_t$ be known at $t$, and let $\E[\varepsilon_t(m_{t+h}-P^\ast_{t+h})]$ vanish as $h$ grows. Then
\[
\E[\varepsilon_t(m_{t+h}-m_t)]\;\longrightarrow\;\E[\varepsilon_t(P^\ast_t-m_t)],
\]
the \omterm{def:mx:decomposing-the-spread:adverse}{adverse-selection component}. The efficient price's own move adds to each trade's impact a term of mean zero and variance $\sigma^2h$, where
$\sigma^2$ is the variance rate of $P^\ast$: over $N$ independent trades it has standard deviation $\sigma\sqrt{h/N}$ in the average.
\end{proposition}

\begin{proof}
Write $\varepsilon_t(m_{t+h}-m_t)=\varepsilon_t(m_{t+h}-P^\ast_{t+h})+\varepsilon_t(P^\ast_{t+h}-P^\ast_t)+\varepsilon_t(P^\ast_t-m_t)$. The middle term has mean
zero: $\varepsilon_t$ is known at $t$ and $P^\ast$'s increments after $t$ have conditional mean zero. The first vanishes by assumption. The middle term
has variance $\E[\varepsilon_t^2(P^\ast_{t+h}-P^\ast_t)^2]=\sigma^2h$ since $\varepsilon_t^2=1$, and the average of $N$ independent copies has variance
$\sigma^2h/N$.
\end{proof}

The proposition is a trade-off. A short horizon has not let the mid catch up, and reads information as spread earned; a long horizon lets the efficient
price's own moves swamp the average. In a real market trades also overlap: a trade's impact at 5 minutes contains the impact of every trade in those
5 minutes, including the rest of its own metaorder.

\Cref{fig:mx:decomposing-the-spread:horizon} shows both effects on the simulated day. The true adverse part is 0.449 ticks, 84\% of the effective
half-spread: the \omterm{def:mx:why-there-is-a-spread:informed}{informed traders}, 21\% of the volume, trade when the efficient price is 2.05 ticks beyond the mid on their side, \omterm{def:mx:why-there-is-a-spread:informed}{noise traders} when it
is 0.025 ticks beyond. The impact climbs from 0.03 ticks at 50 milliseconds to 0.08 at one second, 0.21 at five, 0.35 at fifteen, 0.40 at thirty and 0.48
at one minute, then falls to 0.40 at two minutes and 0.20 at five. The fall is noise, not reversal: the standard error, from 10-minute blocks, grows from
0.07 ticks at one minute to 0.15 at five, a ratio of 2.28 against $\sqrt5=2.24$. The simulator can remove the efficient price's own move,
$\varepsilon_t(P^\ast_{t+h}-P^\ast_t)$, from each trade's impact; the rest settles at 0.46 ticks from one minute on (0.48 at two minutes, 0.46 at five,
standard errors 0.04), 85\% of the effective half-spread, close to the truth. The realised spread falls from 0.50 ticks at 50 milliseconds to 0.056 at
one minute: the providers keep a tenth of what they charge.

\begin{omfigure}
\begin{tikzpicture}
\begin{axis}[width=10cm,height=5.4cm,xmode=log,xlabel={\small horizon $h$ (seconds)},ylabel={\small ticks},tick label style={font=\footnotesize},
  xmin=0.035,xmax=420,ymin=-0.05,ymax=0.62,grid=major,grid style={gray!25},legend columns=2,legend style={font=\footnotesize,at={(0.5,-0.3)},
  anchor=north,draw=none,fill=none,/tikz/every even column/.append style={column sep=8pt}},table/col sep=comma]
\addplot[omDef,very thick,mark=*,error bars/.cd,y dir=both,y explicit] table[x=h,y=impact,y error=se]{figdata/microstructure/05-decomposing-the-spread/horizon.csv};
\addplot[omThm,very thick,mark=square*,error bars/.cd,y dir=both,y explicit] table[x=h,y=clean,y error=clean_se]{figdata/microstructure/05-decomposing-the-spread/horizon.csv};
\addplot[omProp,thick,mark=triangle*] table[x=h,y=realised]{figdata/microstructure/05-decomposing-the-spread/horizon.csv};
\addplot[black,dashed,domain=0.035:420] {0.4491};
\legend{impact,impact net of $P^\ast$'s move,realised half-spread,true adverse part}
\end{axis}
\end{tikzpicture}

\omcaption{Price impact and realised half-spread by horizon on the simulated day (share-weighted, with standard errors from 10-minute blocks); the effective
half-spread is 0.535 ticks. The impact reaches the true adverse part at about one minute; beyond, the efficient price's own moves make it noisy. Data:
\texttt{mx\_decomp.by\_horizon}.}\label{fig:mx:decomposing-the-spread:horizon}
\end{omfigure}

Splitting the trades by the simulator's flag shows who pays. At one minute, informed trades have an effective half-spread of 0.54 ticks, an impact of
2.11 and a realised spread of $-1.57$: each informed share costs the provider a tick and a half. Noise trades pay 0.53, move the mid 0.045 and leave 0.49.
The provider's small margin is the difference between two large numbers.

\section{Structural decompositions}

Before intraday quote data were common, spreads were decomposed from transaction prices alone, with a model of how prices respond to trades.

Glosten and Harris (1988) wrote the transaction price as $p_t=m_t+\varepsilon_t(c_0+c_1v_t)$, a transitory part linear in the trade size $v_t$, and
the mid's revision as $m_t-m_{t-1}=\varepsilon_t(z_0+z_1v_t)+y_t$, a permanent part; price changes are then linear in $\Delta\varepsilon_t$,
$\Delta(\varepsilon_tv_t)$, $\varepsilon_t$ and $\varepsilon_tv_t$.

\begin{definition}[Huang--Stoll model]\label{def:mx:decomposing-the-spread:hs}
The \emph{Huang--Stoll model}\index{Huang--Stoll model} (basic form) writes the change of the transaction price between trades as
\[
\Delta p_t=\frac S2\,\Delta\varepsilon_t+\lambda\,\frac S2\,\varepsilon_{t-1}+e_t,
\]
where $S$ is the traded spread and $\lambda$ the share of the half-spread by which the quotes move after a trade: adverse selection plus inventory, which
the basic form cannot separate. Its extensions separate them with a model of the autocorrelation of trade signs.
\end{definition}

\begin{definition}[Madhavan--Richardson--Roomans model]\label{def:mx:decomposing-the-spread:mrr}
The \emph{Madhavan--Richardson--Roomans model}\index{Madhavan--Richardson--Roomans model} (MRR) lets the efficient price move with the surprise in
the order flow and the transaction price bounce around it:
\[
\mu_t=\mu_{t-1}+\theta\left(x_t-\E[x_t\mid x_{t-1}]\right)+u_t,\qquad p_t=\mu_t+\phi\,x_t+\xi_t,
\]
with $x_t$ the trade sign, $\E[x_t\mid x_{t-1}]=\rho x_{t-1}$, $\theta$ the information content of a trade and $\phi$ the cost that does not depend on
information. The implied spread is $2(\phi+\theta)$ and the information share $\theta/(\phi+\theta)$.
\end{definition}

Differencing, $\Delta p_t=(\phi+\theta)x_t-(\phi+\rho\theta)x_{t-1}+u_t+\Delta\xi_t$. MRR estimate $(\theta,\phi,\rho)$ by the generalised method of
moments (One Quant Book 4, chapter~11); with $\rho$ taken as the sample autocorrelation of the signs, the moment conditions that make the error orthogonal
to $x_t$ and $x_{t-1}$ are the normal equations of a least-squares regression, which is what the build does.

\omcode{code/firm/spreaddecomp/firm_spreaddecomp.py}{72}{81}{The MRR model: information and cost from a regression on the current and previous signs.}\label{lst:mx:decomposing-the-spread:mrr}

On the simulated day the three models agree with each other and not with the truth. Huang--Stoll returns a traded spread of 1.07 ticks and
$\lambda=0.09$; MRR an information share of 0.13 with $\rho=0.33$; Glosten--Harris an adverse share of 0.10 for a 100-share trade, falling with size
(the simulator's informed orders are slightly smaller, 109 shares on average against 120). All three assume that the quotes have absorbed a trade's
information by the next trade. In the simulated market the mid moves toward the efficient price over tens of seconds, through quote updates between
trades, and a model built on trade-to-trade changes sees only the first revision: it reads the rest as noise and understates adverse selection by a
factor of six to nine.

\section{The vector-autoregression view}

Hasbrouck (1991) let the data choose the dynamics. Take the mid-quote change $r_t$ from just before trade $t$ to just before trade $t+1$ and the trade
sign $x_t$, and fit a vector autoregression (One Quant Book 4, chapter~20):
\[
r_t=\sum_{i=1}^{L}a_ir_{t-i}+\sum_{i=0}^{L}b_ix_{t-i}+v_{1,t},\qquad x_t=\sum_{i=1}^{L}c_ir_{t-i}+\sum_{i=1}^{L}d_ix_{t-i}+v_{2,t}.
\]
The trade enters the quote equation at lag zero: quotes respond to trades within the interval, not the reverse. The cumulative impulse response of the mid
to a unit trade shock, summed until it stops changing, is the permanent impact of a trade: the information it carries, whatever the path the quotes take
to absorb it. Unlike MRR, it lets the absorption take many trades.

It still takes as many trades as the lags allow. On the simulated day, the permanent impact is 0.14 of the effective half-spread with one lag, 0.28 with
five, 0.38 with ten, 0.49 with twenty, 0.56 with fifty and 0.48 with a hundred (\cref{fig:mx:decomposing-the-spread:methods}): it grows while the lags
cover more of the adjustment, then the extra coefficients add noise. It never reaches the truth, for a second reason: the VAR is linear in the signs,
while in the simulator the \omterm{def:mx:why-there-is-a-spread:informed}{informed traders} trade when the gap $P^\ast-m$ is large, a nonlinear rule the linear response averages away.

\begin{omfigure}
\begin{tikzpicture}
\begin{axis}[width=10cm,height=6.2cm,xbar,bar width=7pt,xlabel={\small adverse-selection share of the effective half-spread},
  tick label style={font=\footnotesize},label style={font=\small},xmin=0,xmax=1,ymin=-0.6,ymax=10.6,ytick={0,1,2,3,4,5,6,7,8,9,10},y dir=reverse,
  yticklabels={truth,realised (1 min),Glosten--Harris (100),Huang--Stoll,MRR,VAR 1 lag,VAR 5 lags,VAR 10 lags,VAR 20 lags,VAR 50 lags,VAR 100 lags},
  grid=major,grid style={gray!25},table/col sep=comma]
\addplot[omDef,fill=omDef!40] table[x=share,y=k]{figdata/microstructure/05-decomposing-the-spread/estimators.csv};
\draw[black,thick,dashed] (axis cs:0.839,-0.6) -- (axis cs:0.839,10.6);
\end{axis}
\end{tikzpicture}

\omcaption{The adverse-selection share of the spread by each method on the simulated day, against the truth (dashed) computed from the efficient price.
Only the realised-spread method at one minute, which sees the quotes, comes close. Data: \texttt{mx\_decomp.estimators}, \texttt{var\_by\_lags}.}\label{fig:mx:decomposing-the-spread:methods}
\end{omfigure}

The lesson is not that the structural models are wrong: each is right about the market it assumes. It is that the method must match how quotes absorb
information in the market at hand, and that where a simulator with a known truth exists, the method should be scored on it before its number is used.

\section{Spreads from daily data}

\begin{definition}[Low-frequency spread estimator]\label{def:mx:decomposing-the-spread:lowfreq}
A \emph{low-frequency spread estimator}\index{low-frequency spread estimator} infers the effective spread from daily (or bar) prices: closes, highs and
lows, for markets and periods without intraday quotes. Roll's estimator (One Quant Book 4, chapter~21) uses the autocovariance of price changes;
Corwin and Schultz (2012) the high--low ranges; Abdi and Ranaldo (2017) closes and mid-ranges.
\end{definition}

Corwin and Schultz observe that the high is usually a buy at the ask and the low a sale at the bid, so a day's range contains the spread once and the
volatility once, and the volatility part grows with the window while the spread part does not. Comparing two single-day ranges with the range over both
days separates them. With $H$ and $L$ the log high and low,
\[
\beta=\sum_{j=0}^{1}(H_{t+j}-L_{t+j})^2,\quad \gamma=\bigl(\max(H_t,H_{t+1})-\min(L_t,L_{t+1})\bigr)^2,\quad
\alpha=\frac{\sqrt{2\beta}-\sqrt\beta}{3-2\sqrt2}-\sqrt{\frac{\gamma}{3-2\sqrt2}},
\]
and the relative spread is $S=2(e^\alpha-1)/(1+e^\alpha)$, set to zero when negative. Abdi and Ranaldo use the close $c_t$ and the mid-range
$\eta_t=\tfrac12(H_t+L_t)$: $S^2=4\,\E[(c_t-\eta_t)(c_t-\eta_{t+1})]$.

Both need the spread to be a visible part of the range. On the simulated day the effective spread is 1.07 basis points of the price. On bars of 30
seconds, one minute, five minutes and fifteen minutes (761, 390, 78 and 26 bars), Corwin--Schultz returns 0.5, 0.6, 2.0 and 2.6 basis points and Abdi--Ranaldo
0.1, 0.0, 3.3 and 0.0: nowhere reliable, because in a liquid market with a one-tick spread most of the range is volatility and the estimators subtract
two large numbers. They were built for daily data on less liquid stocks, whose spreads are a larger share of the range, and are best checked, as here,
against a period where the effective spread is known.

\section{Tutorial: who paid the two cents}\label{tut:mx:decomposing-the-spread:tut}

\begin{tutorial}
\textbf{Goal.} Decompose the simulated day's spread by every method and score each against the truth.
\textbf{End state:} \cref{fig:mx:decomposing-the-spread:horizon,fig:mx:decomposing-the-spread:methods} and the numbers of sections 2 to 5.
\begin{tutsteps}
\item \textbf{The day.} \texttt{mx\_decomp.day()}: the 6.5-hour session of \texttt{firm.tape}, with its efficient price \texttt{v} and the \texttt{informed}
flag of every trade.
\item \textbf{Truth.} \texttt{truth(tape)}: the effective half-spread and $\E[\varepsilon(P^\ast-m)]$, overall and by trader type.
\item \textbf{Horizons.} \texttt{by\_horizon(tape)}: \texttt{decompose} at eight horizons, with block standard errors and the impact net of $P^\ast$'s move.
\item \textbf{Models.} \texttt{estimators(tape)}: Glosten--Harris, Huang--Stoll, MRR and the VAR; \texttt{var\_by\_lags(tape)} for the lag scan.
\item \textbf{Daily data.} \texttt{low\_frequency(tape, bar\_s)} for bars of 30 seconds to 15 minutes; draw with \texttt{fig\_decomp.py}.
\end{tutsteps}
\textbf{What to change next.} Make the liquidity providers of \texttt{firm.tape} slower to follow the efficient price (a lower rate of quote updates)
and watch the horizon at which the impact settles move out; run the decomposition on \texttt{firm.exchsim}'s recorded day, where Book~7's tape agents
trade against the \omterm{def:mx:the-limit-order-book:engine}{matching engine}.
\end{tutorial}

\section{Build: the spread decomposition}\label{bld:mx:decomposing-the-spread:spreaddecomp}

\begin{build}
\bfield{Purpose} The spread measures and decompositions used by the transaction-cost analysis of chapter~19, by the market-quality metrics of chapter~24
and by Book~11's mark-out studies.
\bfield{Interface} \texttt{mid\_at}, \texttt{quoted\_spread}, \texttt{decompose(trades, times, mids, h)}; \texttt{glosten\_harris}, \texttt{huang\_stoll},
\texttt{mrr}; \texttt{hasbrouck\_var(r, x, lags, steps)}; \texttt{corwin\_schultz}, \texttt{abdi\_ranaldo} (Roll's estimator is in
\texttt{firm.spreadmodels}).
\bfield{Rules} Trades carry the aggressor's sign; the mid is the one in force just before the trade; averages are share-weighted; the structural models
are fitted by least squares with numpy; negative low-frequency estimates are set to zero, as the original papers do.
\bfield{Acceptance tests} \texttt{code/firm/spreaddecomp/tests/}: effective equals realised plus impact on hand trades; MRR, Huang--Stoll and
Glosten--Harris recover the parameters of markets simulated from their own models; the VAR recovers a planted permanent impact; Corwin--Schultz and
Abdi--Ranaldo recover a planted spread of 1\% from minute prices.
\bfield{Stretch} The three-way Huang--Stoll decomposition with sign autocorrelation; MRR by full GMM with standard errors; Hasbrouck's information share
across venues.
\end{build}

\begin{omsources}
\item L.~R.~Glosten and L.~E.~Harris, ``Estimating the components of the bid/ask spread'', \emph{Journal of Financial Economics} 21(1), 1988.
\item R.~D.~Huang and H.~R.~Stoll, ``The components of the bid-ask spread: a general approach'', \emph{Review of Financial Studies} 10(4), 1997.
\item A.~Madhavan, M.~Richardson and M.~Roomans, ``Why do security prices change? A transaction-level analysis of NYSE stocks'', \emph{Review of Financial
Studies} 10(4), 1997.
\item J.~Hasbrouck, ``Measuring the information content of stock trades'', \emph{Journal of Finance} 46(1), 1991.
\item S.~A.~Corwin and P.~Schultz, ``A simple way to estimate bid-ask spreads from daily high and low prices'', \emph{Journal of Finance} 67(2), 2012.
\item F.~Abdi and A.~Ranaldo, ``A simple estimation of bid-ask spreads from daily close, high, and low prices'', \emph{Review of Financial Studies} 30(12), 2017.
\item U.S. Securities and Exchange Commission, \emph{Disclosure of Order Execution Information}, Release 34-99679, 89 FR 26428, 2024; exemptive order,
Release 34-105136, 2026.
\end{omsources}

\section{Exercises}

\begin{exercise}[$\star$]\label{exo:mx:decomposing-the-spread:1}
The market is 100.00 bid, 100.04 offered. A buy of 1,000 shares executes at 100.03, and a minute later the mid is 100.025. Give the quoted and effective
half-spreads, the price improvement, the realised half-spread and the impact per share, and the effective spread in basis points.
\end{exercise}

\begin{exercise}[$\star$]\label{exo:mx:decomposing-the-spread:2}
A sale executes at 49.98 with the mid at 50.00; five minutes later the mid is 49.97. Give the effective and realised half-spreads and the impact, and say
what a negative realised spread means for the liquidity provider.
\end{exercise}

\begin{exercise}[$\star$]\label{exo:mx:decomposing-the-spread:3}
The basic \omterm{def:mx:decomposing-the-spread:hs}{Huang--Stoll model} returns a traded spread of 0.04 and $\lambda=0.6$. By how much do the quotes move after a buy, and how much of the half-spread
is left to cover processing costs?
\end{exercise}

\begin{exercise}[$\star\star$]\label{exo:mx:decomposing-the-spread:4}
An MRR fit gives $\theta=0.012$, $\phi=0.008$ and $\rho=0.3$. Give the implied spread and information share, and the expected price change for a buy that
follows a buy and for a buy that follows a sale. Why do they differ?
\end{exercise}

\begin{exercise}[$\star\star$]\label{exo:mx:decomposing-the-spread:5}
On the simulated day the impact is 0.48 ticks at one minute and 0.20 at five, with standard errors 0.07 and 0.15. Was the one-minute move reversed? What
does the ratio of the standard errors tell you, and what does the impact net of $P^\ast$'s move show?
\end{exercise}

\begin{exercise}[$\star\star$]\label{exo:mx:decomposing-the-spread:6}
Two consecutive days have highs 100.10 and 100.12 and lows 99.94 and 99.96. Compute $\beta$, $\gamma$ and the Corwin--Schultz spread.
\end{exercise}

\begin{exercise}[$\star\star\star$]\label{exo:mx:decomposing-the-spread:7}
\emph{Coding.} On the simulated day, split the trades by the \texttt{informed} flag and decompose each group at one minute. Who pays the provider's margin,
and how does the share-weighted mix give back the overall realised spread of 0.056 ticks?
\end{exercise}

\begin{exercise}[$\star\star\star$]\label{exo:mx:decomposing-the-spread:8}
\emph{Find the flaw.} ``Our algorithm's child orders show a negative five-minute realised spread for the market makers who fill them, so our flow is
informed and we deserve better prices.''
\end{exercise}

\section{Problem: Who Paid the Two Cents?}

\begin{problem}[{Weekend problem --- who paid the two cents?}]\label{pb:mx:decomposing-the-spread:1}
A liquidity provider wants to know how much of her spread pays for being picked off, and a regulator wants to know which horizon to report. On the
simulated market the truth is known; score every method against it.

\medskip\noindent\textbf{Part I --- Spreads and impact.}
\begin{enumerate}
\item Write the identity between effective, realised and impact half-spreads and prove it.
\item What are the \omterm{def:mx:decomposing-the-spread:quoted}{quoted spread} and the share-weighted effective half-spread of the simulated day?
\item Define the \omterm{def:mx:decomposing-the-spread:adverse}{adverse-selection component} and compute the truth: in ticks and as a share.
\item How far is the efficient price from the mid when informed and \omterm{def:mx:why-there-is-a-spread:informed}{noise traders} trade?
\item Give the impact at 50 milliseconds, one second, fifteen seconds and one minute.
\end{enumerate}

\medskip\noindent\textbf{Part II --- The horizon.}
\begin{enumerate}[resume]
\item State and prove the convergence proposition.
\item Why does the impact fall between one and five minutes?
\item What is the impact net of the efficient price's move, and why can only a simulation compute it?
\item Give the realised spread at 50 milliseconds and at one minute, and say what the provider keeps.
\item Which of the five Rule 605 horizons is closest to where the impact settles here?
\end{enumerate}

\medskip\noindent\textbf{Part III --- Models.}
\begin{enumerate}[resume]
\item Derive the MRR price-change equation.
\item What do Huang--Stoll, MRR and Glosten--Harris return?
\item Why do they understate adverse selection on this market?
\item What does the VAR return with 1, 5, 20, 50 and 100 lags, and why is the pattern not monotonic?
\item Why can a linear VAR not reach the truth here?
\end{enumerate}

\medskip\noindent\textbf{Part IV --- Daily data and the verdict.}
\begin{enumerate}[resume]
\item Explain the idea of the Corwin--Schultz estimator.
\item What do Corwin--Schultz and Abdi--Ranaldo return on 1-minute and 5-minute bars, against the effective spread?
\item When can low-frequency estimators be trusted?
\item State the \emph{named result}: the adverse-selection share estimated by each method against the simulator's truth, and the horizon at which the
realised spread settles.
\item In one sentence: who paid the two cents?
\end{enumerate}
\end{problem}

\section{Interview questions}

\begin{interviewq}[$\star$ \iqroles{trader, researcher}]\label{iq:mx:decomposing-the-spread:1}
Define the effective and realised spreads. What does the difference between them measure?
\end{interviewq}

\begin{interviewq}[$\star\star$ \iqroles{researcher}]\label{iq:mx:decomposing-the-spread:2}
How do you choose the horizon of a realised spread, and what goes wrong if it is too short or too long?
\end{interviewq}

\begin{interviewq}[$\star\star$ \iqroles{researcher}]\label{iq:mx:decomposing-the-spread:3}
Write the MRR model and explain how its two parameters are identified.
\end{interviewq}

\begin{interviewq}[$\star\star$ \iqroles{trader}]\label{iq:mx:decomposing-the-spread:4}
Your market-making desk earns a positive one-second realised spread but loses money. How is that possible?
\end{interviewq}

\begin{interviewq}[$\star\star$ \iqroles{researcher}]\label{iq:mx:decomposing-the-spread:5}
What is Hasbrouck's VAR measuring, and why does the trade enter the quote equation at lag zero?
\end{interviewq}

\begin{interviewq}[$\star\star\star$ \iqroles{researcher}]\label{iq:mx:decomposing-the-spread:6}
You need spreads for a sample of stocks in the 1970s with only daily data. What do you use, and how do you check it?
\end{interviewq}
